\chapter{Introduction}
\label{ch:introduction}

\section{Purpose and Scope}

The SIRC-1 CPU is sold on its own, so it is documented in three manuals that stand apart from any system built around
it. SIRCADA and SIRCAS are made for the SIRCUS system and each has a manual of its own. This manual completes the set:
it describes the system that the parts are built into. The six manuals have these scopes:

\begin{table}[H]
	\centering
	\begin{tabular}{lp{8.2cm}}
		\toprule
		\textbf{Manual} & \textbf{Scope}                                                                                               \\
		\midrule
		\archmanual{}   & ISA, instruction encodings, canonical assembly notation, exceptions, programmer-visible state                \\
		\asmmanual{}    & Complete assembly syntax, directives, symbols, meta-instructions, and toolchain usage                        \\
		\hwmanual{}     & Pinout, bus interface, timing, electrical characteristics                                                    \\
		\videomanual{}  & The display adapter: package and pinout, memory map, registers, data formats, video modes, timing            \\
		\audiomanual{}  & The sound processor: package and pinout, memory map, registers, data formats, voices, filter, sequencer      \\
		\sysmanual{}    & The console: clocks, bus, address map, interrupts, reset, connections and timing between chips (this manual) \\
		\bottomrule
	\end{tabular}
	\caption{The SIRC-1 and SIRCUS Manual Set}
	\label{tab:manual-set}
\end{table}

Where a chip's own manual is authoritative for something, this manual says so and names the chapter. Where the
connection of two chips creates a fact that neither manual could state alone, such as the number of audio samples in a
video frame, this manual is authoritative.

\section{The Parts of the System}

SIRCUS has three chips that matter to a program, each on the one 16-bit bus, and the memories and circuits around them.

\begin{table}[H]
	\centering
	\small
	\begin{tabular}{llp{5.2cm}}
		\toprule
		\textbf{Part}     & \textbf{Package}            & \textbf{Function}                                                                               \\
		\midrule
		SIRC-1 CPU        & 64-pin DIP                  & Runs the program; the only master of the bus                                                    \\
		SIRCADA           & 52-pin PLCC                 & Display adapter: tiles, sprites and the video signal; holds its own video memory                \\
		SIRCAS            & Not yet specified           & Sound processor: six voices, filter and sequencers; has its own audio memory                    \\
		Audio memory      & Not yet specified           & Static memory of 32,768 words, used by SIRCAS alone                                             \\
		Converter         & TDA1541A, 28-pin DIP        & Philips dual 16-bit converter: turns SIRCAS's serial output into two analogue currents          \\
		Output amplifiers & Not yet specified           & Turn the converter's currents into voltages (a reference circuit uses two NE5534 or equivalent) \\
		Working memory    & Not yet specified           & 128 KB (65,536 words) of RAM for the program's data and stack                                   \\
		System ROM        & Not yet specified           & Holds the start-up code and the driver routines that programs call                              \\
		Cartridge         & Connector not yet specified & Holds a game's program and data in a ROM; its first words are the reset vector                  \\
		Serial I/O        & MC68681 DUART               & Motorola dual serial port: links two consoles, a modem or a debugger                            \\
		Select logic      & 74LS138, 16-pin DIP         & Decodes A23 to A21 into a chip select for each device; a latch holds the boot overlay           \\
		Clock circuit     & Not yet specified           & Makes the master clock and divides it by two to make \texttt{CLKI}                              \\
		\bottomrule
	\end{tabular}
	\caption{The Parts of SIRCUS}
	\label{tab:parts}
\end{table}

The video memory is inside SIRCADA, and the CPU reaches it through the device's own segment (\videomanual{}, memory map
chapter). The audio memory is outside SIRCAS, but it is not on the CPU's bus: the CPU reaches it only through a window
of SIRCAS's segment, and SIRCAS makes every access to it itself (\audiomanual{}, memory map chapter). The CPU's bus
therefore carries traffic for only six things: the working memory, the system ROM, the cartridge, SIRCADA, SIRCAS and
the serial device.

\section{Block Diagram}

Figure~\ref{fig:block-diagram} shows the connections. One clock rail feeds all three chips, one bus joins them, and
each of the other two chips has outputs that the CPU never sees directly: the picture, and the sound.

\begin{figure}[H]
	\centering
	\begin{tikzpicture}[
			scale=0.68, transform shape,
			chip/.style={draw, thick, fill=gray!8, align=center, font=\bfseries\small, minimum width=2.7cm, minimum height=1.6cm},
			box/.style={draw, rounded corners=2pt, fill=sirclightgray, align=center, font=\small, minimum width=2.5cm, minimum height=0.95cm},
			flow/.style={-{Stealth[length=2mm]}, thick},
			rail/.style={very thick},
			note/.style={font=\scriptsize, text=sircgray, align=center}
		]
		% Clock
		\node[box] (osc) at (0,5.6) {Master clock oscillator};
		\node[box] (div) at (0,4.0) {Divide by 2};
		\draw[flow] (osc) -- (div);
		\draw[rail] (-5,2.9) -- (5,2.9);
		\draw[thick] (div.south) -- (0,2.9);
		\node[note, anchor=south west] at (0.2,3.0) {\texttt{CLKI}};
		% Chips
		\node[chip] (cpu) at (-5,0) {SIRC-1\\CPU};
		\node[chip] (ppu) at (0,0) {SIRCADA};
		\node[chip] (snd) at (5,0) {SIRCAS};
		\foreach \c in {cpu,ppu,snd} { \draw[flow] (\c |- 0,2.9) -- (\c.north); }
		% Bus
		\draw[rail] (-7.7,-2.0) -- (6.4,-2.0);
		\foreach \c in {cpu,ppu,snd} { \draw[thick, <->] (\c.south) -- (\c |- 0,-2.0); }
		\node[note, anchor=south west, align=left] at (0.55,-1.97) {A0--A23, D0--D15\\\activelow{BAS}, BRW, \activelow{BACK}};
		% Below the bus
		\node[box, minimum width=2.0cm] (ram) at (-6.6,-4.2) {Working\\RAM};
		\node[box, minimum width=2.0cm] (rom) at (-4.2,-4.2) {System\\ROM};
		\node[box, minimum width=2.0cm] (cart) at (-1.8,-4.2) {Cartridge};
		\node[box, minimum width=2.2cm] (sel) at (4.4,-4.2) {Select\\logic};
		\node[box, minimum width=2.0cm] (ser) at (6.9,-4.2) {Serial\\I/O};
		\node[box, minimum width=2.0cm] (sconn) at (9.2,-4.2) {Serial\\connector};
		\draw[thick, <->] (ser.north) -- (ser |- 0,-2.0);
		\draw[flow, <->] (ser) -- (sconn);
		\foreach \c in {ram,rom,cart,sel} { \draw[thick, <->] (\c.north) -- (\c |- 0,-2.0); }
		% Outputs
		\node[box] (disp) at (1.4,-4.9) {Display\\R, G, B, \activelow{CSYNC}};
		\draw[flow] (0.6,-0.8) -- (0.6,-4.43);
		% Audio side
		\node[box] (amem) at (9.2,0.9) {Audio\\memory};
		\node[box] (dac) at (9.2,-0.9) {Converter};
		\draw[thick, <->] (snd.east |- amem) -- (amem.west);
		\draw[flow] (snd.east |- dac) -- (dac.west);
		\node[box] (out) at (9.2,-2.9) {Stereo\\audio out};
		\draw[flow] (dac) -- (out);
	\end{tikzpicture}
	\caption{SIRCUS Block Diagram}
	\label{fig:block-diagram}
\end{figure}

The interrupt and reset connections are not drawn; Chapter~\ref{ch:interrupts} gives them, and
Chapter~\ref{ch:connections} lists every connection of every chip.

\section{Conventions}

\begin{itemize}
	\item \textbf{Addresses} are word addresses of 24 bits, as in the \hwmanual{}. A segment is 65,536 words, and the
	      eight high address bits select it.
	\item \textbf{Active-low signals} carry a bar over the name and are asserted when low.
	\item \textbf{Clock cycles} are cycles of \texttt{CLKI}, the clock every chip takes, unless a figure says otherwise.
	      The master clock is twice as fast.
	\item \textbf{Both regions.} Figures are given for an NTSC console and a PAL console. A figure without a region
	      holds for both.
	\item \textbf{Not yet specified.} A value that the design has not fixed is marked ``not yet specified''. It is not
	      undefined, and no value is implied.
\end{itemize}
